\chapter{Conclusion and Outlook}
\label{ch:conclusion}

\draftnote{this chapter is a placeholder skeleton --- rewrite once
\cref{ch:discussion} is filled in with concrete numbers, since the claims
below should restate quantified results, not generic statements.}

\section{Summary}

This thesis set out to determine whether sub-wavelength displacement and
material change can be recovered from time-lapse GPR surveys once classical
amplitude-based migration imaging has stopped resolving them.
\Cref{ch:resolution} quantified that amplitude floor for two stationary
point scatterers across three independent migration algorithms.
\Cref{ch:tl-horizontal,ch:tl-vertical} showed that the same floor limits
amplitude \emph{differencing} of a single displaced scatterer, both
laterally and vertically, with and without additive noise. A 2D Fourier
phase-plane shift estimator was then derived (\cref{ch:theory}) and
validated (\cref{ch:phase-plane}) directly on the same datasets, and applied
to a more field-realistic, spatially distributed fluid-front scenario
(\cref{ch:fluidflow}), where the same fit was shown to decouple a purely
geometric displacement from a purely dielectric material change using a
single weighted least-squares plane fit.

\section{Outlook for Field Application}

\draftnote{expand with concrete next steps. Candidates already implied by
the methodology and theory chapters: (1) validate the pipeline on a real
zero-offset field survey with a known, controlled sub-wavelength
displacement (e.g.\ a target on a calibrated micrometre stage) to test the
field pre-processing assumptions of \cref{sec:meth-phaseplane} (dewow,
time-zero correction, trace re-binning) that the synthetic data in this
thesis sidesteps by construction; (2) run the three outstanding robustness
experiments flagged in \cref{ch:discussion}; (3) extend the fluid-front
model of \cref{ch:fluidflow} to a swept range of fracture thickness and
fluid contrast, and to a front geometry that is not perfectly straight; (4)
quantify the velocity-recalibration procedure needed when soil or ice
moisture genuinely changes between baseline and monitor surveys, since
\cref{sec:meth-phaseplane} noted that this is otherwise indistinguishable
from a true vertical shift $\Dz$.}
